\documentclass[11pt,a4paper]{article}
\usepackage{curso}
\begin{document}
\portada{Ejercicios --- Lección}{10.4}{Transformers dispersos}
\section{Ejercicios}
\begin{enumerate}[label=\textbf{E\arabic*.}]
\item (B) Ancho con grado 8 y presupuesto de 8192 conexiones.
\item (B) ¿Por qué sube el error con grado muy bajo?
\item (B) Capas para cubrir 4096 entradas con grado 16: patrón local frente a mundo pequeño.
\item (B) ¿Por qué las capas dispersas pueden ser lentas en una GPU?
\item (I) Con grado 16, ¿cuántas conexiones tiene una capa de ancho 4096? ¿Y la densa equivalente?
\item (I) Proponga un patrón de mundo pequeño de grado 16: cuántas conexiones locales y cuántos saltos largos.
\item (A) Repita el barrido iso-FLOP con otro maestro y otro presupuesto; ¿se mueve el óptimo?
\end{enumerate}
\section{Soluciones}
\begin{solucion}
\textbf{E1.} $8192/8=1024$.

\textbf{E2.} Cada unidad depende de muy pocas entradas y no puede representar las combinaciones que el maestro necesita.

\textbf{E3.} Local: $4096/16=256$; mundo pequeño: $\log_{16}4096=3$.

\textbf{E4.} Acceso irregular a memoria y poca reutilización; las unidades tensoriales están optimizadas para bloques densos.

\textbf{E5.} $4096\cdot16=65\,536$ frente a $4096^2\approx1{,}7\times10^7$: 256 veces menos.

\textbf{E6.} Por ejemplo 12 locales y 4 saltos aleatorios.

\textbf{E7.} Suele moverse con el presupuesto y la tarea.
\end{solucion}
\end{document}
